% Options for packages loaded elsewhere
% Options for packages loaded elsewhere
\PassOptionsToPackage{unicode}{hyperref}
\PassOptionsToPackage{hyphens}{url}
\PassOptionsToPackage{dvipsnames,svgnames,x11names}{xcolor}
%
\documentclass[
  12pt,
  letterpaper,
]{article}
\usepackage{xcolor}
\usepackage[margin=1in]{geometry}
\usepackage{amsmath,amssymb}
\setcounter{secnumdepth}{-\maxdimen} % remove section numbering
\usepackage{iftex}
\ifPDFTeX
  \usepackage[T1]{fontenc}
  \usepackage[utf8]{inputenc}
  \usepackage{textcomp} % provide euro and other symbols
\else % if luatex or xetex
  \usepackage{unicode-math} % this also loads fontspec
  \defaultfontfeatures{Scale=MatchLowercase}
  \defaultfontfeatures[\rmfamily]{Ligatures=TeX,Scale=1}
\fi
\usepackage{lmodern}
\ifPDFTeX\else
  % xetex/luatex font selection
\fi
% Use upquote if available, for straight quotes in verbatim environments
\IfFileExists{upquote.sty}{\usepackage{upquote}}{}
\IfFileExists{microtype.sty}{% use microtype if available
  \usepackage[]{microtype}
  \UseMicrotypeSet[protrusion]{basicmath} % disable protrusion for tt fonts
}{}
\makeatletter
\@ifundefined{KOMAClassName}{% if non-KOMA class
  \IfFileExists{parskip.sty}{%
    \usepackage{parskip}
  }{% else
    \setlength{\parindent}{0pt}
    \setlength{\parskip}{6pt plus 2pt minus 1pt}}
}{% if KOMA class
  \KOMAoptions{parskip=half}}
\makeatother


\usepackage{longtable,booktabs,array}
\usepackage{calc} % for calculating minipage widths
% Correct order of tables after \paragraph or \subparagraph
\usepackage{etoolbox}
\makeatletter
\patchcmd\longtable{\par}{\if@noskipsec\mbox{}\fi\par}{}{}
\makeatother
% Allow footnotes in longtable head/foot
\IfFileExists{footnotehyper.sty}{\usepackage{footnotehyper}}{\usepackage{footnote}}
\makesavenoteenv{longtable}
\usepackage{graphicx}
\makeatletter
\newsavebox\pandoc@box
\newcommand*\pandocbounded[1]{% scales image to fit in text height/width
  \sbox\pandoc@box{#1}%
  \Gscale@div\@tempa{\textheight}{\dimexpr\ht\pandoc@box+\dp\pandoc@box\relax}%
  \Gscale@div\@tempb{\linewidth}{\wd\pandoc@box}%
  \ifdim\@tempb\p@<\@tempa\p@\let\@tempa\@tempb\fi% select the smaller of both
  \ifdim\@tempa\p@<\p@\scalebox{\@tempa}{\usebox\pandoc@box}%
  \else\usebox{\pandoc@box}%
  \fi%
}
% Set default figure placement to htbp
\def\fps@figure{htbp}
\makeatother


% definitions for citeproc citations
\NewDocumentCommand\citeproctext{}{}
\NewDocumentCommand\citeproc{mm}{%
  \begingroup\def\citeproctext{#2}\cite{#1}\endgroup}
\makeatletter
 % allow citations to break across lines
 \let\@cite@ofmt\@firstofone
 % avoid brackets around text for \cite:
 \def\@biblabel#1{}
 \def\@cite#1#2{{#1\if@tempswa , #2\fi}}
\makeatother
\newlength{\cslhangindent}
\setlength{\cslhangindent}{1.5em}
\newlength{\csllabelwidth}
\setlength{\csllabelwidth}{3em}
\newenvironment{CSLReferences}[2] % #1 hanging-indent, #2 entry-spacing
 {\begin{list}{}{%
  \setlength{\itemindent}{0pt}
  \setlength{\leftmargin}{0pt}
  \setlength{\parsep}{0pt}
  % turn on hanging indent if param 1 is 1
  \ifodd #1
   \setlength{\leftmargin}{\cslhangindent}
   \setlength{\itemindent}{-1\cslhangindent}
  \fi
  % set entry spacing
  \setlength{\itemsep}{#2\baselineskip}}}
 {\end{list}}
\usepackage{calc}
\newcommand{\CSLBlock}[1]{\hfill\break\parbox[t]{\linewidth}{\strut\ignorespaces#1\strut}}
\newcommand{\CSLLeftMargin}[1]{\parbox[t]{\csllabelwidth}{\strut#1\strut}}
\newcommand{\CSLRightInline}[1]{\parbox[t]{\linewidth - \csllabelwidth}{\strut#1\strut}}
\newcommand{\CSLIndent}[1]{\hspace{\cslhangindent}#1}



\setlength{\emergencystretch}{3em} % prevent overfull lines

\providecommand{\tightlist}{%
  \setlength{\itemsep}{0pt}\setlength{\parskip}{0pt}}



 


% -----------------------
% CUSTOM PREAMBLE STUFF
% -----------------------

% -----------------
% Title block stuff
% -----------------

% Abstract
\usepackage[runin]{abstract}
\renewcommand{\abstractnamefont}{\sffamily\small\bfseries}
\renewcommand{\abstracttextfont}{\sffamily\small}
\setlength{\absleftindent}{5pt}
\setlength{\absrightindent}{\absleftindent}

% Title
\usepackage{titling}
\pretitle{\par\begin{flushleft}\LARGE\sffamily\bfseries}
\posttitle{\par\end{flushleft}\vskip 10pt}

% Keywords
\newenvironment{keywords}
{\small\sffamily{\sffamily\small\bfseries{Keywords.}}}

% Authors
\usepackage{orcidlink}  % Create automatic ORCID icons/links
%\renewcommand{\and}{\end{tabular} \hskip 3em \begin{tabular}[t]{@{\hspace{0em}}l@{}}}
\preauthor{\begin{flushleft}
           \lineskip 1.5em}
\postauthor{\end{flushleft}}

% ------------------
% Section headings
% ------------------
\usepackage{titlesec}
\titleformat*{\section}{\Large\sffamily\bfseries\raggedright}
\titleformat*{\subsection}{\large\sffamily\bfseries\raggedright}
\titleformat*{\subsubsection}{\normalsize\sffamily\bfseries\raggedright}
\titleformat*{\paragraph}{\small\sffamily\bfseries\raggedright}

%\titlespacing{<command>}{<left>}{<before-sep>}{<after-sep>}
% Starred version removes indentation in following paragraph
\titlespacing*{\section}{0em}{2em}{0.1em}
\titlespacing*{\subsection}{0em}{1.25em}{0.1em}
\titlespacing*{\subsubsection}{0em}{0.75em}{0em}

% ------------------
% Headers/Footers
% ------------------
\usepackage{fancyhdr}
\pagestyle{fancy}
\fancyhf{}
\fancyhead[L,C,R]{}
\fancyfoot[L,C]{}
\fancyfoot[R]{\thepage}
\renewcommand{\headrulewidth}{1pt}
\fancypagestyle{plain}{%
    \renewcommand{\headrulewidth}{0pt}%
    \fancyhf{}%
    \fancyfoot[R]{\thepage}%
}
\renewcommand\footnoterule{\rule{\linewidth}{0.1pt}\vspace{5pt}}

% ------------------
% Captions
% ------------------
\usepackage[labelfont=bf,labelsep=period]{caption}
\captionsetup[figure]{font=footnotesize,justification=raggedright,singlelinecheck=false,format=hang}


% ---------------------------
% END CUSTOM PREAMBLE STUFF
% ---------------------------
\makeatletter
\@ifpackageloaded{caption}{}{\usepackage{caption}}
\AtBeginDocument{%
\ifdefined\contentsname
  \renewcommand*\contentsname{Table of contents}
\else
  \newcommand\contentsname{Table of contents}
\fi
\ifdefined\listfigurename
  \renewcommand*\listfigurename{List of Figures}
\else
  \newcommand\listfigurename{List of Figures}
\fi
\ifdefined\listtablename
  \renewcommand*\listtablename{List of Tables}
\else
  \newcommand\listtablename{List of Tables}
\fi
\ifdefined\figurename
  \renewcommand*\figurename{Figure}
\else
  \newcommand\figurename{Figure}
\fi
\ifdefined\tablename
  \renewcommand*\tablename{Table}
\else
  \newcommand\tablename{Table}
\fi
}
\@ifpackageloaded{float}{}{\usepackage{float}}
\floatstyle{ruled}
\@ifundefined{c@chapter}{\newfloat{codelisting}{h}{lop}}{\newfloat{codelisting}{h}{lop}[chapter]}
\floatname{codelisting}{Listing}
\newcommand*\listoflistings{\listof{codelisting}{List of Listings}}
\makeatother
\makeatletter
\makeatother
\makeatletter
\@ifpackageloaded{caption}{}{\usepackage{caption}}
\@ifpackageloaded{subcaption}{}{\usepackage{subcaption}}
\makeatother
\usepackage{bookmark}
\IfFileExists{xurl.sty}{\usepackage{xurl}}{} % add URL line breaks if available
\urlstyle{same}
\hypersetup{
  pdftitle={State Renewable Energy Policies and Renewable Electricity Generation: An Analysis of Renewable Portfolio Standards in 2008--2009},
  pdfauthor={Samane Afshari},
  pdfkeywords={Renewable Portfolio Standards, renewable energy
production, US States},
  colorlinks=true,
  linkcolor={blue},
  filecolor={Maroon},
  citecolor={Blue},
  urlcolor={red},
  pdfcreator={LaTeX via pandoc}}



\title{State Renewable Energy Policies and Renewable Electricity
Generation: An Analysis of Renewable Portfolio Standards in 2008--2009}
% subtitles do not seem to work with article class?
%%\subtitle{}

\author{
{\bfseries \normalsize Samane Afshari~\orcidlink{0000-0001-7237-8131}}%
\thanks{Corresponding author.} \\%
 \small University of Oregon,  \\%
{\footnotesize \url{samaa@uoregon.edu}} \\\vspace{10pt}
}

\predate{}
\postdate{}
\date{}
\begin{document}

% for some reason this does not work in header
\renewcommand{\abstractname}{Abstract.}

% add the short title to the fancy header
\fancyhead[R]{Paper}
\fancyhead[L]{Samane Afshari}

\maketitle
%\noindent \rule{\linewidth}{.5pt}
\begin{abstract}
\ldots..
\end{abstract}
\begin{keywords}
\def\sep{;\ }
Renewable Portfolio Standards\sep renewable energy production\sep 
US States
\end{keywords}
%\noindent \rule{\linewidth}{.5pt}


\section{Introduction}\label{introduction}

The transition toward renewable energy has become a central policy
objective in the United States as concerns over climate change, carbon
emissions, and fossil fuel dependence continue to intensify. In the
absence of a unified federal leadership, state governments have
increasingly emerged as major regulator in renewable energy governance
through different energy-related policy instruments.

Renewable Portfolio Standards (RPS) are among the most widely used
state-level policy mechanisms because they require utilities to generate
a specified share of electricity from renewable sources such as wind,
solar, geothermal, and biomass. By 2008, more than half of U.S. states
had adopted mandatory RPS programs, reflecting a broader shift toward
decentralized energy governance and state-led climate action.

Despite the rapid expansion of state renewable energy policies, there
remains ongoing debate regarding whether these policies actually
increase renewable energy production. Using PRS and Renewable
Electricity Generation as the main variables, this study examines
whether state-level renewable energy policies increase renewable
electricity generation across U.S. states during 2008 and 2009.

Existing literature has produced mixed findings on the effectiveness of
Renewable Portfolio Standards. Early studies often relied on qualitative
analyses or case studies that highlighted successful examples of
renewable energy deployment but provided limited empirical evidence
regarding broader policy effectiveness. More recent quantitative
research has attempted to evaluate the relationship between RPS policies
and renewable energy generation. For example, Carley (2009) examined
whether state RPS adoption increased the share of renewable electricity
generation across states and found that while RPS policies increased the
total amount of renewable generation, they did not significantly
increase renewable energy's share within total electricity production.

This study builds directly on Carley's (2009) evaluation of state
renewable electricity policies, as it extends Carley's core policy
question into the 2008--2009 period, when more states had adopted RPS
programs and policy designs had become more developed. This study offers
a focused short-term reassessment of whether state renewable energy
mandates were beginning to shift the composition of state electricity
systems, not simply increase renewable generation in absolute terms.
(Carley 2009).

\section{Background}\label{background}

\ldots.

\section{Data and Methods}\label{data-and-methods}

Independent Variable (IV): State Renewable Portfolio Standards (RPS) The
independent variable is derived from the dataset U.S. State Electricity
Resource Standards: 2024 Status Update, compiled by the Lawrence
Berkeley National Laboratory (LBNL). The dataset provides information on
state-level Renewable Portfolio Standards (RPS) and Clean Energy
Standards (CES) across U.S. states, including policy adoption dates,
renewable energy targets, and compliance timelines. For this study, the
dataset can be used to measure whether a state has adopted an RPS policy
and the strength of the policy target over time, though the research
only focus on 2008 and 2009.

Dependent Variable (DV): Renewable Electricity Generation Share The
dependent variable is drawn from the U.S. Energy Information
Administration (EIA) Annual Energy Review renewable electricity
generation dataset. The dataset provides state-level information on
electricity generation and the percentage of generation derived from
renewable energy sources in 2008 and 2009. In this study, the dependent
variable is the share of electricity generation from renewable energy
sources, with particular attention to the ``percent non-hydro
renewable'' measure, which captures renewable generation from sources
such as wind, solar, geothermal, and biomass while excluding
conventional hydroelectric power. This measure is appropriate because it
reflects the types of renewable energy sources most directly influenced
by modern RPS policies.

\section{Results}\label{results}

\section{Conclusions}\label{conclusions}

\section{References}\label{references}

\phantomsection\label{refs}
\begin{CSLReferences}{1}{0}
Data source: Lawrence Berkeley National Laboratory RPS Dataset Data
source: U.S. Energy Information Administration Renewable Energy Annual
Data

\phantomsection\label{refs}
\begin{CSLReferences}{1}{0}
\bibitem[\citeproctext]{ref-carley2009}
Carley, Sanya. 2009. {``State Renewable Energy Electricity Policies: An
Empirical Evaluation of Effectiveness.''} \emph{Energy Policy} 37 (8):
3071--81. \url{https://doi.org/10.1016/j.enpol.2009.03.062}.

\end{CSLReferences}

\bibitem[\citeproctext]{ref-carley2009}
Carley, Sanya. 2009. {``State Renewable Energy Electricity Policies: An
Empirical Evaluation of Effectiveness.''} \emph{Energy Policy} 37 (8):
3071--81. \url{https://doi.org/10.1016/j.enpol.2009.03.062}.

\end{CSLReferences}




\end{document}
